\documentclass[11pt]{article}
\usepackage[margin=1in]{geometry}
\usepackage{amsmath,amssymb,amsthm}
\usepackage{mathrsfs}
\usepackage{graphicx}
\usepackage{hyperref}

\newtheorem{theorem}{Theorem}
\newtheorem{lemma}{Lemma}
\newtheorem{corollary}{Corollary}

\newcommand{\B}{\mathscr{B}}
\newcommand{\N}{\mathbb{N}}
\newcommand{\ellone}{\ell_1}

\title{A partial answer to the non-reflexive operator-algebra question in arXiv:1707.08399}
\author{Run fa\_banach\_001, lane 6}
\date{2026-06-29}

\begin{document}
\maketitle

\section*{Source question}

Beanland, Kania and Laustsen ask in Section 3 of
\emph{The algebras of bounded operators on the Tsirelson and Baernstein spaces
are not Grothendieck spaces}, arXiv:1707.08399:

\begin{quote}
Is there a non-reflexive Banach space $X$ for which $\B(X)$ is a
Grothendieck space? If such a space $X$ exists, both $X$ and $X^*$ would
necessarily be Grothendieck spaces, because $\B(X)$ contains complemented
copies of them.
\end{quote}

\begin{center}
\includegraphics[width=0.92\linewidth]{figures/open_problem_crop.png}
\end{center}

The result below rules out all infinite-dimensional spaces of the form
$C(K)$.

\section*{Result}

\begin{theorem}
Let $K$ be an infinite compact Hausdorff space. Then the Banach space
$\B(C(K))$ is not Grothendieck.
\end{theorem}

\begin{corollary}
No infinite-dimensional $C(K)$ space can be a positive example for the
non-reflexive question in arXiv:1707.08399.
\end{corollary}

\section*{Proof}

We use two standard permanence facts, included for completeness.

\begin{lemma}
If a Banach space $Z$ is Grothendieck, then every complemented subspace of
$Z$ is Grothendieck.
\end{lemma}

\begin{proof}
Let $Y$ be complemented in $Z$, and let $P:Z\to Y$ be a bounded projection.
Suppose $(y_n^*)$ is a weak-star null sequence in $Y^*$. Then
$z_n^*=y_n^*\circ P$ is weak-star null in $Z^*$. If $Z$ is Grothendieck, then
$(z_n^*)$ is weakly null in $Z^*$. For $y^{**}\in Y^{**}$ define
$z^{**}\in Z^{**}$ by $z^{**}(z^*)=y^{**}(z^*|_Y)$. Since
$z^{**}(z_n^*)=y^{**}(y_n^*)$, the sequence $(y_n^*)$ is weakly null.
\end{proof}

\begin{lemma}
The space $\ell_1$ is not Grothendieck.
\end{lemma}

\begin{proof}
A separable Grothendieck space is reflexive: if $X$ is separable, the dual
unit ball $B_{X^*}$ is weak-star compact metrizable, so every sequence in it
has a weak-star convergent subsequence. If $X$ is Grothendieck, that
subsequence is weakly convergent. Hence $B_{X^*}$ is weakly sequentially
compact, and the Eberlein--Smulian theorem makes it weakly compact. Thus
$X^*$, and therefore $X$, is reflexive. Since $\ell_1$ is separable and not
reflexive, it is not Grothendieck.
\end{proof}

Now let $K$ be infinite. Choose a countably infinite subset
$D=\{t_1,t_2,\ldots\}$ of $K$. Under the Riesz representation identification
$C(K)^*=M(K)$, the closed subspace
\[
  A_D=\left\{\sum_{j=1}^\infty a_j\delta_{t_j} : (a_j)\in \ell_1\right\}
\]
is isometric to $\ell_1(D)$. It is also $1$-complemented in $M(K)$. Indeed,
for $\mu\in M(K)$ put
\[
  P_D\mu=\sum_{j=1}^\infty \mu(\{t_j\})\,\delta_{t_j}.
\]
The total variation of this atomic part is bounded by $\|\mu\|$, so
$P_D$ is a contractive projection from $M(K)$ onto $A_D$.

Since $A_D$ is isometric to $\ell_1$ and $\ell_1$ is not Grothendieck, the
first permanence lemma shows that $M(K)=C(K)^*$ is not Grothendieck.

Finally, $C(K)^*$ is complemented in $\B(C(K))$ by rank-one operators. Fix
$t_0\in K$, let $\mathbf{1}$ denote the constant-one function, and define
\[
  J:M(K)\to \B(C(K)),\qquad
  J(\mu)(f)=\mu(f)\mathbf{1}.
\]
Define also
\[
  Q:\B(C(K))\to M(K),\qquad
  Q(T)(f)=\delta_{t_0}(T f).
\]
Then $QJ=\mathrm{id}_{M(K)}$, so $J(M(K))$ is complemented in $\B(C(K))$.
If $\B(C(K))$ were Grothendieck, then $M(K)$ would be Grothendieck by the
first lemma, contradicting the preceding paragraph. Therefore $\B(C(K))$ is
not Grothendieck.

\section*{Verification and limitations}

The proof uses only complemented-subspace permanence of the Grothendieck
property, the elementary fact that separable Grothendieck spaces are
reflexive, the Riesz representation theorem for $C(K)^*$, and explicit
rank-one embeddings. No computational verification is required.

This is a partial result. It excludes the classical non-reflexive family
$C(K)$, $K$ infinite, but does not solve the source's general non-reflexive
problem. In particular, it does not rule out a nonseparable Banach space $X$
for which both $X$ and $X^*$ are Grothendieck, and it does not address the
source's questions about $\B(\ell_p)$, weak Hilbert spaces, super-reflexive
spaces, or reflexive spaces with unconditional FDDs.

\section*{Novelty check}

The local run indexes and existing attempts for arXiv:1707.08399 were checked
on 2026-06-29; no existing packet recorded this $C(K)$ exclusion. Bounded web
searches for the exact operator-algebra questions around $\B(\ell_p)$ and
Grothendieck operator algebras did not locate a later exact resolution of the
source's open-problem list. Because the present argument is short and uses
standard facts, it should be reviewed as a likely-valid partial exclusion with
possible folklore status rather than as a claimed full solution.

\begin{thebibliography}{9}
\bibitem{BKL}
K. Beanland, T. Kania and N. J. Laustsen,
\emph{The algebras of bounded operators on the Tsirelson and Baernstein spaces
are not Grothendieck spaces}, arXiv:1707.08399.
\end{thebibliography}

\end{document}
